% GENERATED FILE. Edit canonical lessons, then run npm run generate:books.
% canonical-source-hash: fnv1a64:17af11265207c0d7
% canonical-lessons: TA-C388-opputal, TA-C388-ankikari, TA-C388-niyami, TA-C388-terivi, TA-C388-amarvu

\chapter{Approval and Appointments}
\label{ch:ta-approval-and-appointments}
\hlchaptermodality{388}

\begin{chapteropening}
\textbf{By the end of this chapter:} \emph{I can say that something was approved and someone appointed.}

Complete the last lesson of chapter 388: I can say that something was approved and someone appointed.
\end{chapteropening}

\section[\texorpdfstring{\ta{ஒப்புதல்} (opputal)}{ஒப்புதல் (opputal)}]{\ta{ஒப்புதல்} (opputal) — approval; consent}

\label{lesson:TA-C388-opputal}



\begin{quote}
Before the new one: say the Tamil for a pledge, then the Tamil for to fulfil.
\end{quote}

\subsection*{You'll want to know: \ta{ஒப்புதல்}}
\textbf{\ta{ஒப்புதல்}} — \emph{opputal} — \textquotedblleft{}approval; consent\textquotedblright{}.

\begin{quote}
\ta{இந்தத்} \ta{திட்டத்துக்கு} \ta{அரசின்} \ta{ஒப்புதல்} \ta{கிடைத்தது}.
\end{quote}

\begin{quote}
\emph{intat tiṭṭattukku araciṉ opputal kiṭaittatu.}
\end{quote}

\begin{quote}
This plan got the approval of the government.
\end{quote}

\ta{ஒப்புதல்} \ta{கொடு}, to approve.

\begin{cousinweb}
\textbf{\ta{ஒப்புதல்}} (\emph{opputal}) is native: the verbal noun in -\ta{தல்} of \emph{oppu}, \textquotedblleft{}to agree, accept\textquotedblright{}, from \emph{o}, \textquotedblleft{}to be suited, acceptable\textquotedblright{} (DEDR 924).
\end{cousinweb}

\begin{grammarlens}[title={What you've built}]
One more word for saying what happens if.
\end{grammarlens}

\subsection*{Your turn}
\begin{itemize}
\raggedright
  \item \emph{Say it:} \emph{opputal}
  \item \emph{Say it:} \emph{opputal}, once more
  \item \emph{Say it:} \emph{intat tiṭṭattukku araciṉ opputal kiṭaittatu.} — This plan got the approval of the government.
  \item \emph{Recall:} say the Tamil for a pledge, then the Tamil for to fulfil, then say \emph{opputal} again
\end{itemize}

\begin{tcolorbox}[breakable,colback=teal!4,colframe=teal!35!black,title={Before you move on}]
What does \ta{ஒப்புதல்} mean? (\textquotedblleft{}Approval; consent\textquotedblright{}.) Say it once more.
\end{tcolorbox}
\section[\texorpdfstring{\ta{அங்கீகரி} (aṅkīkari)}{அங்கீகரி (aṅkīkari)}]{\ta{அங்கீகரி} (aṅkīkari) — to recognise officially; to accredit}

\label{lesson:TA-C388-ankikari}



\begin{quote}
Before the new one: say the Tamil for to fulfil, then the Tamil for approval.
\end{quote}

\subsection*{You'll want to know: \ta{அங்கீகரி}}
\textbf{\ta{அங்கீகரி}} — \emph{aṅkīkari} — \textquotedblleft{}to recognise officially; to accredit\textquotedblright{}.

\begin{quote}
\ta{அரசு} \ta{இந்தப்} \ta{பள்ளியை} \ta{அங்கீகரித்தது}.
\end{quote}

\begin{quote}
\emph{aracu intap paḷḷiyai aṅkīkarittatu.}
\end{quote}

\begin{quote}
The government recognised this school.
\end{quote}

Past (I …): \ta{அங்கீகரித்தேன்}, \emph{aṅkīkarittēṉ}.

What is recognised takes -\ta{ஐ}. The noun is \ta{அங்கீகாரம்}, recognition.

\begin{cousinweb}
\textbf{\ta{அங்கீகரி}} (\emph{aṅkīkari}) is a verb built on Sanskrit \emph{aṅgīkāra}, \textquotedblleft{}agreement, promise\textquotedblright{} (Monier-Williams), from \emph{aṅgī-kṛ}, \textquotedblleft{}to make one's own, to accept\textquotedblright{}.
\end{cousinweb}

\begin{grammarlens}[title={What you've built}]
One more word for saying what happens if.
\end{grammarlens}

\subsection*{Your turn}
\begin{itemize}
\raggedright
  \item \emph{Say it:} \emph{aṅkīkari}
  \item \emph{Say it:} \emph{aṅkīkari}, once more
  \item \emph{Say it:} \emph{aracu intap paḷḷiyai aṅkīkarittatu.} — The government recognised this school.
  \item \emph{Recall:} say the Tamil for to fulfil, then the Tamil for approval, then say \emph{aṅkīkari} again
\end{itemize}

\begin{tcolorbox}[breakable,colback=teal!4,colframe=teal!35!black,title={Before you move on}]
What does \ta{அங்கீகரி} mean? (\textquotedblleft{}To recognise officially; to accredit\textquotedblright{}.) Say it once more.
\end{tcolorbox}
\section[\texorpdfstring{\ta{நியமி} (niyami)}{நியமி (niyami)}]{\ta{நியமி} (niyami) — to appoint}

\label{lesson:TA-C388-niyami}



\begin{quote}
Before the new one: say the Tamil for approval, then the Tamil for to recognise officially.
\end{quote}

\subsection*{You'll want to know: \ta{நியமி}}
\textbf{\ta{நியமி}} — \emph{niyami} — \textquotedblleft{}to appoint\textquotedblright{}.

\begin{quote}
\ta{நிறுவனம்} \ta{ஒரு} \ta{புதிய} \ta{மேலாளரை} \ta{நியமித்தது}.
\end{quote}

\begin{quote}
\emph{niṟuvaṉam oru putiya mēlāḷarai niyamittatu.}
\end{quote}

\begin{quote}
The company appointed a new manager.
\end{quote}

Past (I …): \ta{நியமித்தேன்}, \emph{niyamittēṉ}.

Who is appointed takes -\ta{ஐ}; as what: \ta{மேலாளராக} \ta{நியமி}.

\begin{cousinweb}
\textbf{\ta{நியமி}} (\emph{niyami}) is a verb built on Sanskrit \emph{niyama}, \textquotedblleft{}restraint; fixed law\textquotedblright{} (CDIAL 7250), with the loan-verb ending -\ta{இ}: to fix someone in a post.
\end{cousinweb}

\begin{grammarlens}[title={What you've built}]
One more word for saying what happens if.
\end{grammarlens}

\subsection*{Your turn}
\begin{itemize}
\raggedright
  \item \emph{Say it:} \emph{niyami}
  \item \emph{Say it:} \emph{niyami}, once more
  \item \emph{Say it:} \emph{niṟuvaṉam oru putiya mēlāḷarai niyamittatu.} — The company appointed a new manager.
  \item \emph{Recall:} say the Tamil for approval, then the Tamil for to recognise officially, then say \emph{niyami} again
\end{itemize}

\begin{tcolorbox}[breakable,colback=teal!4,colframe=teal!35!black,title={Before you move on}]
What does \ta{நியமி} mean? (\textquotedblleft{}To appoint\textquotedblright{}.) Say it once more.
\end{tcolorbox}
\section[\texorpdfstring{\ta{தெரிவி} (terivi)}{தெரிவி (terivi)}]{\ta{தெரிவி} (terivi) — to inform; to convey}

\label{lesson:TA-C388-terivi}



\begin{quote}
Before the new one: say the Tamil for to recognise officially, then the Tamil for to appoint.
\end{quote}

\subsection*{You'll want to know: \ta{தெரிவி}}
\textbf{\ta{தெரிவி}} — \emph{terivi} — \textquotedblleft{}to inform; to convey\textquotedblright{}.

\begin{quote}
\ta{அவர்} \ta{தன்} \ta{முடிவை} \ta{எங்களுக்குத்} \ta{தெரிவித்தார்}.
\end{quote}

\begin{quote}
\emph{avar taṉ muṭivai eṅkaḷukkut terivittār.}
\end{quote}

\begin{quote}
He informed us of his decision.
\end{quote}

Past (I …): \ta{தெரிவித்தேன்}, \emph{terivittēṉ}.

The person told takes -\ta{க்கு}: \ta{கருத்துத்} \ta{தெரிவி}, to give one's view.

\begin{cousinweb}
\textbf{\ta{தெரிவி}} (\emph{terivi}) is native Dravidian: DEDR 3419 lists Tamil \emph{teri}, \textquotedblleft{}to be seen, understood\textquotedblright{}, and its causative \emph{terivi}, \textquotedblleft{}to make known\textquotedblright{}.
\end{cousinweb}

\begin{grammarlens}[title={What you've built}]
One more word for saying what happens if.
\end{grammarlens}

\subsection*{Your turn}
\begin{itemize}
\raggedright
  \item \emph{Say it:} \emph{terivi}
  \item \emph{Say it:} \emph{terivi}, once more
  \item \emph{Say it:} \emph{avar taṉ muṭivai eṅkaḷukkut terivittār.} — He informed us of his decision.
  \item \emph{Recall:} say the Tamil for to recognise officially, then the Tamil for to appoint, then say \emph{terivi} again
\end{itemize}

\begin{tcolorbox}[breakable,colback=teal!4,colframe=teal!35!black,title={Before you move on}]
What does \ta{தெரிவி} mean? (\textquotedblleft{}To inform; to convey\textquotedblright{}.) Say it once more.
\end{tcolorbox}
\section[\texorpdfstring{\ta{அமர்வு} (amarvu)}{அமர்வு (amarvu)}]{\ta{அமர்வு} (amarvu) — a session; a sitting}

\label{lesson:TA-C388-amarvu}



\begin{quote}
Before the new one: say the Tamil for to appoint, then the Tamil for to inform.
\end{quote}

\subsection*{You'll want to know: \ta{அமர்வு}}
\textbf{\ta{அமர்வு}} — \emph{amarvu} — \textquotedblleft{}a session; a sitting\textquotedblright{}.

\begin{quote}
\ta{நாடாளுமன்றத்தின்} \ta{அடுத்த} \ta{அமர்வு} \ta{திங்கள்கிழமை} \ta{தொடங்கும்}.
\end{quote}

\begin{quote}
\emph{nāṭāḷumaṉṟattiṉ aṭutta amarvu tiṅkaḷkiḻamai toṭaṅkum.}
\end{quote}

\begin{quote}
The next session of parliament begins on Monday.
\end{quote}

From \ta{அமர்}, to sit.

\begin{cousinweb}
\textbf{\ta{அமர்வு}} (\emph{amarvu}) is native Dravidian: DEDR 161 lists Tamil \emph{amar}, \textquotedblleft{}to abide, remain, rest\textquotedblright{}, and \emph{amarvu}, \textquotedblleft{}abode\textquotedblright{}; a session is a sitting together.
\end{cousinweb}

\begin{grammarlens}[title={What you've built}]
One more word for saying what happens if.
\end{grammarlens}

\subsection*{Your turn}
\begin{itemize}
\raggedright
  \item \emph{Say it:} \emph{amarvu}
  \item \emph{Say it:} \emph{amarvu}, once more
  \item \emph{Say it:} \emph{nāṭāḷumaṉṟattiṉ aṭutta amarvu tiṅkaḷkiḻamai toṭaṅkum.} — The next session of parliament begins on Monday.
  \item \emph{Recall:} say the Tamil for to appoint, then the Tamil for to inform, then say \emph{amarvu} again
\end{itemize}

\begin{tcolorbox}[breakable,colback=teal!4,colframe=teal!35!black,title={Before you move on}]
What does \ta{அமர்வு} mean? (\textquotedblleft{}A session; a sitting\textquotedblright{}.) Say it once more.
\end{tcolorbox}
